\dictchar{अँ / अं}

\bhojentry{अँकवार}{\K{\ba\bchandra\bka\bva\bmaa\bra}}{Añkwaar}{हि. स्त्री. | ठेठ भोजपुरी}{\freqcommon}%
{आलिंगन, छाती से लगाना, गोद।}%
{अँकवार (Añkwaar) / कोरवा}%
{Embrace, open-armed hug, lap.}
\bhojexample{लइका माई के अँकवार में छिपा गइल।}{The child hid in his mother's embrace.}

\bhojentry{अँकवारना}{\K{\ba\bchandra\bka\bva\bmaa\bra\bna\bmaa}}{Añkwaarna}{हि. क्रि. स. | ठेठ भोजपुरी}{\freqcommon}%
{गले लगाना, छाती से चिपटाना।}%
{अँकवार लिहल (Añkwaar Lihal) / छाती से लगावल}%
{To embrace, clasp to chest.}

\bhojentry{अँकवारी}{\K{\ba\bchandra\bka\bva\bmaa\bra\bmii}}{Añkwari}{हि. स्त्री. | ठेठ भोजपुरी}{\freqcommon}%
{गोद, क्रोड़, अंक।}%
{कोरा (Kora) / अँकवारी}%
{Lap, open arms.}

\bhojentry{अँकाना}{\K{\ba\bchandra\bka\bmaa\bna\bmaa}}{Añkana}{हि. क्रि. स. | ठेठ भोजपुरी}{\freqcommon}%
{मूल्य या भार का अनुमान लगवाना, कूतवाना।}%
{अंटकावल (Añtkawal) / कूतवावल}%
{To cause to be evaluated, appraised.}

\bhojentry{अँकुरी}{\K{\ba\bchandra\bka\bmu\bra\bmii}}{Añkuri}{हि. स्त्री. | ठेठ भोजपुरी/कृषि}{\freqcommon}%
{पौधे की बेल का मुड़ा भाग, अंकुर, कांटी।}%
{अँकुरी (Añkuri) / काँटा}%
{Tendril, sprout shoot, small hook.}

\bhojentry{अँखुआ}{\K{\ba\bchandra\bkha\bmu\baa}}{Añkhua}{हि. पुं. | ठेठ भोजपुरी}{\freqcommon}%
{अंकुर, बीज से फूटता हुआ नया पौधा।}%
{अँखुआ (Añkhua) / कल्ला}%
{Sprout, germinating seed shoot.}

\bhojentry{अँखुआना}{\K{\ba\bchandra\bkha\bmu\baa\bna\bmaa}}{Añkhuana}{हि. क्रि. | ठेठ भोजपुरी}{\freqcommon}%
{अंकुरित होना, अंखुआ निकलना।}%
{अँखुआइल (Añkhuail) / कल्ला निकलब}%
{To sprout, germinate.}

\bhojentry{अँगलू}{\K{\ba\bchandra\bga\bla\bmuu}}{Añgaluu}{हि. पुं. | ठेठ भोजपुरी}{\freqformal}%
{दर्जी का उंगली में पहनने का छल्ला (अंगुष्ठाना)।}%
{अँगलू (Añgaluu) / अंगुष्ठाना}%
{Tailor's thimble, finger guard.}

\bhojentry{अँगुरी}{\K{\ba\bchandra\bga\bmu\bra\bmii}}{Añguri}{हि. स्त्री. | ठेठ भोजपुरी}{\freqcommon}%
{हाथ या पैर की उंगली।}%
{अँगुरी (Añguri) / आँगुर}%
{Finger, toe.}

\bhojentry{अँगुरी काटना}{\K{\ba\bchandra\bga\bmu\bra\bmii} \ \K{\bka\bmaa\btta\bna\bmaa}}{Añguri Kaatna}{हि. मुहा. | ठेठ भोजपुरी}{\freqcommon}%
{आश्चर्य या पश्चाताप में दांतों तले उंगली दबाना।}%
{अँगुरी काटल (Añguri kaatal)}%
{To bite fingers in amazement/regret.}

\bhojentry{अँगूठा}{\K{\ba\bchandra\bga\bmuu\bttha\bmaa}}{AñguuTha}{हि. पुं. | तद्भव}{\freqcommon}%
{हाथ-पैर की पहली बड़ी उंगली, अंगुष्ठ।}%
{अँगूठा (AñguuTha) / अंगुठा}%
{Thumb, big toe.}

\bhojentry{अँगूठी}{\K{\ba\bchandra\bga\bmuu\bttha\bmii}}{AñguuThii}{हि. स्त्री. | तद्भव}{\freqcommon}%
{उंगली की मुद्रिका, छल्ला।}%
{अँगूठी (AñguuThii) / छाप}%
{Finger ring.}

\bhojentry{अँगूठा दिखाना}{\K{\ba\bchandra\bga\bmuu\bttha\bmaa} \ \K{\bda\bmi\bkha\bmaa\bna\bmaa}}{AñguuTha Dikhana}{हि. मुहा. | ठेठ भोजपुरी}{\freqcommon}%
{साफ मना कर देना, धोखा देना।}%
{अँगूठा देखावल (AñguuTha dekhaawal)}%
{To refuse flatly, renege.}

\bhojentry{अँगौछा}{\K{\ba\bchandra\bga\bmau\bcha\bmaa}}{Añgaucha}{हि. पुं. | ठेठ भोजपुरी}{\freqcommon}%
{सूती देह पोंछने का तौलिया, गमछा।}%
{अँगौछा (Añgaucha) / गमछा}%
{Traditional cotton towel.}

\bhojentry{अँगार / अँगारा}{\K{\ba\bchandra\bga\bmaa\bra}}{Añgaar}{हि. पुं. | तद्भव}{\freqcommon}%
{दहकता कोयला या लकड़ी।}%
{अँगार (Añgaar) / दहकत कोइला}%
{Burning coal, embers.}

\bhojentry{अँगारिन}{\K{\ba\bchandra\bga\bmaa\bra\bmi\bna}}{Añgaarin}{हि. स्त्री. | ठेठ भोजपुरी}{\freqrare}%
{क्रोधित स्त्री, चिंगारी जैसी।}%
{अँगारिन (Añgaarin) / गुस्सैल मेहरारू}%
{Fiery/hot-tempered woman.}

\bhojentry{अँगेइया}{\K{\ba\bchandra\bga\bme\bai\bya\bmaa}}{Añgeiya}{हि. स्त्री. | ठेठ भोजपुरी}{\freqcommon}%
{भोज-भात या दावत का निमंत्रण।}%
{अँगेइया (Añgeiya) / न्योता}%
{Feast invitation.}

\bhojentry{अँगना}{\K{\ba\bchandra\bga\bna\bmaa}}{Añgna}{हि. पुं. | ठेठ भोजपुरी}{\freqcommon}%
{घर के भीतर का खुला चौकोर स्थान।}%
{अँगना (Añgna) / आँगना}%
{Courtyard, open yard.}

\bhojentry{अँचाना}{\K{\ba\bchandra\bca\bmaa\bna\bmaa}}{Añcaana}{हि. क्रि. स. | ठेठ भोजपुरी}{\freqcommon}%
{भोजन के बाद हाथ-मुंह धोना।}%
{अँचाना (Añcaana) / हाथ-मुंह धोवल}%
{To wash hands and mouth after meal.}

\bhojentry{अँचवना}{\K{\ba\bchandra\bca\bva\bna\bmaa}}{Añcavna}{हि. पुं. | ठेठ भोजपुरी}{\freqformal}%
{हथेली में लिया गया आचमन का पानी।}%
{अँचवना (Añcavna) / आचमन के पानी}%
{Water taken in palm for rinsing mouth.}

\bhojentry{अँचरा}{\K{\ba\bchandra\bca\bra\bmaa}}{Añcra}{हि. पुं. | ठेठ भोजपुरी}{\freqcommon}%
{साड़ी का पल्लू या आंचल।}%
{अँचरा (Añcra) / पल्लू}%
{Sari pallu, skirt border.}

\bhojentry{अँछल}{\K{\ba\bchandra\bcha\bla}}{Añchal}{हि. वि. | ठेठ भोजपुरी}{\freqrare}%
{बिना छाना हुआ गंदा या निर्मल जल।}%
{अँछल पानी (Añchal paani)}%
{Unfiltered or raw water.}

\bhojentry{अँझुराना}{\K{\ba\bchandra\bjha\bmu\bra\bmaa\bna\bmaa}}{Añjhuraana}{हि. क्रि. | ठेठ भोजपुरी}{\freqcommon}%
{उलझना, विवाद में फंसना।}%
{अँझुराइल (Añjhuraail)}%
{To get entangled, caught in dispute.}

\bhojentry{अँझुराइल}{\K{\ba\bchandra\bjha\bmuu\bra\bmaa\bai\bla}}{Añjhuraail}{हि. वि. | ठेठ भोजपुरी}{\freqcommon}%
{उलझा हुआ, विवादित।}%
{अँझुराइल (Añjhuraail)}%
{Entangled, complicated.}

\bhojentry{अँझुरी}{\K{\ba\bchandra\bjha\bmu\bra\bmii}}{Añjhuri}{हि. स्त्री. | ठेठ भोजपुरी}{\freqcommon}%
{दोनों हथेलियों का सम्पुट।}%
{अँझुरी (Añjhuri) / अंजलि}%
{Cupped palms.}

\bhojentry{अँटइया}{\K{\ba\bchandra\btta\bai\bya\bmaa}}{Añtaiya}{हि. स्त्री. | ठेठ भोजपुरी}{\freqcommon}%
{कपड़े या रस्सी की छोटी गांठ।}%
{अँटइया (Añtaiya) / गांठ}%
{Small knot, bundle.}

\bhojentry{अँटकाना}{\K{\ba\bchandra\btta\bka\bmaa\bna\bmaa}}{Añtkaana}{हि. क्रि. स. | ठेठ भोजपुरी}{\freqcommon}%
{फंसाना, बीच में रोकना।}%
{अँटकावल (Añtkawal)}%
{To cause to get stuck/halted.}

\bhojentry{अँटिया}{\K{\ba\bchandra\btta\bmi\bya\bmaa}}{Añtiya}{हि. स्त्री. | ठेठ भोजपुरी/कृषि}{\freqcommon}%
{फसल या घास की छोटी पूली।}%
{अँटिया (Añtiya) / आँटी}%
{Small crop sheaf.}

\bhojentry{अँटी}{\K{\ba\bchandra\btta\bmii}}{Añtii}{हि. स्त्री. | ठेठ भोजपुरी}{\freqcommon}%
{सूत या धागे की छोटी लच्छी।}%
{अँटी (Añtii) / लच्छी}%
{Small skein/spool of thread.}

\bhojentry{अँडवा}{\K{\ba\bchandra\bdda\bva\bmaa}}{Añdwa}{हि. पुं. | ठेठ भोजपुरी}{\freqcommon}%
{अंडी का बीज।}%
{अँडवा (Añdwa) / रेड़ी के बीज}%
{Castor seed.}

\bhojentry{अँडुवा}{\K{\ba\bchandra\bdda\bmu\bva\bmaa}}{Añduwa}{हि. पुं. | ठेठ भोजपुरी/पशु}{\freqcommon}%
{बिना बधाई (असबध) किया गया बैल या सांड।}%
{अँडुवा बैल (Añduwa bail)}%
{Uncastrated bull.}

\bhojentry{अँतराना}{\K{\ba\bchandra\bta\bra\bmaa\bna\bmaa}}{Añtraana}{हि. क्रि. | ठेठ भोजपुरी}{\freqcommon}%
{दूरी बनाना, एक दिन बीच करना।}%
{अँतरा के (Añtra ke)}%
{To keep distance, skip days.}

\bhojentry{अँतरी}{\K{\ba\bchandra\bta\bra\bmii}}{Añtri}{हि. स्त्री. | तद्भव}{\freqcommon}%
{आंत, आंत्र।}%
{अँतरी (Añtri) / आंत}%
{Intestine, entrails.}

\bhojentry{अँदाज़}{\K{\ba\bchandra\bda\bmaa\bja}}{Añdaaz}{हि. पुं. | फारसी आगत}{\freqcommon}%
{अनुमान, ढंग, हाव-भाव।}%
{अँदाज़ (Añdaaz) / अंदाज}%
{Manner, style, estimate.}

\bhojentry{अँधियारा}{\K{\ba\bchandra\bdha\bmi\bya\bmaa\bra\bmaa}}{Añdhiyaara}{हि. पुं. | तद्भव}{\freqcommon}%
{अंधकार, अंधेरा।}%
{अँधियारा (Añdhiyaara) / अंहार}%
{Darkness, gloom.}

\bhojentry{अँधियारी}{\K{\ba\bchandra\bdha\bmi\bya\bmaa\bra\bmii}}{Añdhiyaarii}{हि. स्त्री. | ठेठ भोजपुरी}{\freqcommon}%
{अंधेरी रात, आंखों की पट्टी।}%
{अँधियारी (Añdhiyaarii)}%
{Dark night, blindfold.}

\bhojentry{अँधार}{\K{\ba\bchandra\bdha\bmaa\bra}}{Añdhaar}{हि. पुं. | ठेठ भोजपुरी}{\freqcommon}%
{घोर अंधकार।}%
{अँधार (Añdhaar) / घोर अंहार}%
{Deep darkness.}

\bhojentry{अँधूप}{\K{\ba\bchandra\bdha\bmuu\bpa}}{Añdhuup}{हि. वि. | ठेठ भोजपुरी}{\freqrare}%
{अंधभक्त, पूरी तरह अनजान।}%
{अँधूप (Añdhuup) / अंधा भरोसा}%
{Blind follower, ignorant.}

\bhojentry{अँसन}{\K{\ba\bchandra\bsa\bna}}{Añsan}{हि. पुं. | तद्भव}{\freqformal}%
{उपवास, अनशन।}%
{अँसन (Añsan) / अनशन}%
{Fasting, hunger strike.}

\bhojentry{अँसुआ}{\K{\ba\bchandra\bsa\bmu\baa}}{Añsua}{हि. पुं. | ठेठ भोजपुरी}{\freqcommon}%
{आंसू, लोर।}%
{अँसुआ (Añsua) / लोर (Lor)}%
{Tears, weeping drop.}

\bhojentry{अँसुआई}{\K{\ba\bchandra\bsa\bmu\baa\bai}}{Añsuaai}{हि. वि. | ठेठ भोजपुरी}{\freqcommon}%
{आंसुओं से भरी आंखें।}%
{अँसुआई आँखि (Añsuaai aankhi)}%
{Tearful, teary-eyed.}

\bhojentry{अँसुवाए}{\K{\ba\bchandra\bsa\bmu\bva\bmaa\bme}}{Añsuwaaye}{हि. क्रि. | ठेठ भोजपुरी}{\freqcommon}%
{आंखों में आंसू भर आना।}%
{आँखि अँसुवा गइल (Aankhi añsuwa gail)}%
{To well up with tears.}

\bhojentry{अंजोर}{\K{\ba\banus\bja\bmo\bra}}{Anjor}{हि. पुं. | ठेठ भोजपुरी}{\freqcommon}%
{प्रकाश, उजाला, रोशनी।}%
{अंजोर (Anjor) / उजाला}%
{Light, brightness, illumination.}
\bhojexample{भोर में पूरा गाँव में अंजोर हो गइल।}{At dawn the whole village lit up.}

\bhojentry{अंजोरिया}{\K{\ba\banus\bja\bmo\bra\bmi\bya\bmaa}}{Anjoriya}{हि. स्त्री. | ठेठ भोजपुरी}{\freqcommon}%
{चांदनी रात, शुक्ल पक्ष की रात।}%
{अंजोरिया (Anjoriya) / इजोरिया}%
{Moonlit night, bright lunar fortnight.}

\bhojentry{अंजन}{\K{\ba\banus\bja\bna}}{Anjan}{सं. पुं. | तत्सम}{\freqformal}%
{आंखों में लगाने का काजल, सुरुमा।}%
{काजल (Kaajal) / अंजन}%
{Kohl, mascara, eye salve.}

\bhojentry{अंजाम}{\K{\ba\banus\bja\bmaa\bma}}{Anjaam}{हि. पुं. | अरबी/फारसी आगत}{\freqcommon}%
{परिणाम, अंत, नतीजा।}%
{अंजाम (Anjaam) / नतीजा}%
{Result, outcome, conclusion.}

\bhojentry{अंटकल}{\K{\ba\banus\btta\bka\bla}}{Antkal}{हि. पुं. | ठेठ भोजपुरी}{\freqcommon}%
{अनुमान, अटकल, अंदाजा।}%
{अंटकल (Antkal) / अंदाजा}%
{Guess, estimate, conjecture.}

\bhojentry{अंटा}{\K{\ba\banus\btta\bmaa}}{Anta}{हि. पुं. | ठेठ भोजपुरी}{\freqcommon}%
{कांच या पत्थर की गोली, अंटा।}%
{अंटा (Anta) / गोली}%
{Marble, spherical stone.}

\bhojentry{अंडी}{\K{\ba\banus\bdda\bmii}}{Andi}{हि. स्त्री. | ठेठ भोजपुरी}{\freqcommon}%
{रेड़ का पौधा या उसका बीज जिससे तेल बनता है।}%
{अंडी (Andi) / रेड़ी}%
{Castor plant/bean.}

\bhojentry{अंडा}{\K{\ba\banus\bdda\bmaa}}{Anda}{हि. पुं. | तद्भव}{\freqcommon}%
{पक्षियों या जीवों का डिंब।}%
{अंडा (Anda)}%
{Egg.}

\bhojentry{अंडेल}{\K{\ba\banus\bdda\bme\bla}}{Andel}{हि. वि. | ठेठ भोजपुरी}{\freqcommon}%
{अंडा देने वाली चिड़िया या मुर्गी।}%
{अंडेल (Andel)}%
{Egg-laying bird or hen.}

\bhojentry{अंतर}{\K{\ba\banus\bta\bra}}{Antar}{सं. पुं. | तद्भव}{\freqcommon}%
{भेद, दूरी, भिन्नता।}%
{अंतर (Antar) / भेद}%
{Difference, distance, gap.}

\bhojentry{अंतरंग}{\K{\ba\banus\bta\bra\banus\bga}}{Antrang}{सं. वि. | तत्सम}{\freqformal}%
{घनिष्ठ, आंतरिक।}%
{घनिष्ठ (Ghanishth)}%
{Intimate, internal.}

\bhojentry{अंतर्राष्ट्रीय}{\K{\ba\banus\bta\bra\bvir\bra\bmaa\bvir\bssa\btta\bvir\bra\bmii\bya}}{Antarraashtriya}{सं. वि. | तत्सम}{\freqcommon}%
{विभिन्न राष्ट्रों के बीच।}%
{देशन के बीच}%
{International.}

\bhojentry{अंतरिक्ष}{\K{\ba\banus\bta\bra\bmi\bka\bvir\bsha}}{Antariksh}{सं. पुं. | तत्सम}{\freqcommon}%
{आकाश, खगोल।}%
{आसमान / खगोल}%
{Outer space.}

\bhojentry{अंतरिक्ष यान}{\K{\ba\banus\bta\bra\bmi\bka\bvir\bsha} \ \K{\bya\bmaa\bna}}{Antariksh Yaan}{सं. पुं. | तत्सम}{\freqcommon}%
{अंतरिक्ष में जाने वाला वाहन।}%
{स्पेसक्राफ्ट}%
{Spacecraft.}

\bhojentry{अंतराल}{\K{\ba\banus\bta\bra\bmaa\bla}}{Antraal}{सं. पुं. | तत्सम}{\freqformal}%
{बीच का स्थान या समय।}%
{बीच के समय}%
{Interval, gap.}

\bhojentry{अंतिम}{\K{\ba\banus\bta\bmi\bma}}{Antim}{सं. वि. | तत्सम}{\freqcommon}%
{आखिरी, सबसे अंत का।}%
{अंतिम (Antim) / आखिरी}%
{Final, last, ultimate.}

\bhojentry{अंधा}{\K{\ba\banus\bdha\bmaa}}{Andha}{हि. पुं. / वि. | तद्भव}{\freqcommon}%
{दृष्टिहीन व्यक्ति, अन्हरा।}%
{अन्हरा (Anhra) / अंधा}%
{Blind person.}

\bhojentry{अंधियारा}{\K{\ba\banus\bdha\bmi\bya\bmaa\bra\bmaa}}{Andhiyaara}{हि. पुं. | तद्भव}{\freqcommon}%
{अंधकार, अंहार।}%
{अँधियारा (Añdhiyaara)}%
{Darkness, gloom.}

\bhojentry{अंधेर}{\K{\ba\banus\bdha\bme\bra}}{Andher}{हि. पुं. | तद्भव}{\freqcommon}%
{अन्याय, अव्यवस्था।}%
{अंधेर (Andher) / अन्याय}%
{Anarchy, tyranny, chaos.}

\bhojentry{अंधेर नगरी}{\K{\ba\banus\bdha\bme\bra} \ \K{\bna\bga\bra\bmii}}{Andher Nagari}{हि. स्त्री. | लोकप्रसिद्ध}{\freqcommon}%
{अराजक और अन्यायी राज्य।}%
{अंधेर नगरी}%
{Realm of chaos and injustice.}

\bhojentry{अंधविश्वास}{\K{\ba\banus\bdha\bva\bmi\bsha\bvir\bva\bmaa\bsa}}{Andhvisvaas}{सं. पुं. | तद्भव}{\freqcommon}%
{अंधा भरोसा, कुरीति।}%
{अंधविश्वास}%
{Superstition, blind faith.}

\bhojentry{अंश}{\K{\ba\banus\bsha}}{Ansh}{सं. पुं. | तत्सम}{\freqcommon}%
{भाग, हिस्सा, डिग्री।}%
{अंस (Ans) / हिस्सा}%
{Portion, degree, fraction.}

\bhojentry{अंसार}{\K{\ba\banus\bsa\bmaa\bra}}{Ansaar}{हि. पुं. | अरबी आगत}{\freqformal}%
{मददगार, साथी।}%
{अंसार (Ansaar)}%
{Allies, helpers.}
